\documentclass[../main]{subfiles}
\begin{document}
\chapter{Inducción Magnética}
\img{images/01/magnetic_cover}{Inducción Magnética}
\info{
Contenido}{
  Campo magnético creado por una corriente eléctrica, intensidad de
  campo, inducción,permeabilidad magnética,flujo magnético.
}
\actividad{Resolver los siguientes problemas.}
\begin{enumerate}
  \item Expresar en $Gauss$ una inducción magnética de $B=1,2  Tesla$.
  \item Si la inducción magnética es igual a 18.000 Gauss ¿Cuál será
    su valor en Tesla?
  \item Un solenoide de $l=40cm$ de longitud y $N=1800$ espiras está
    arrollado sobre un núcleo de madera y circulan por el una
    corriente de $I=10A$. Calcular la inducción magnética $B$ en el
    interior del solenoide sabiendo que la permeabilidad magnética de
    la madera es igual a la del aire $\mu_0=4x\pi \times 10^-7
    \frac{T \cdot m}{A}$.
  \item Sobre un anillo de madera cuyo diámetro medio es de
    $d_m=10cm$ se arrolla un devanado de $N=400$ vueltas. Calcular la
    inducción magnética $B$ en el punto de la circunferencia media
    del anillo, si la intensidad de corriente del devanado es de $I=0,5A$.
  \item Calcular la inducción magnética $B$ en el eje de de una
    bobina de $N=400$ espiras devanadas sobre un carrete de cartón de
    $l=25cm$ de longitud y diámetro mucho menor a la misma. La
    intensidad que circula es de $I=2A$.
  \item Un solenoide de $N=500$ espiras está constituido por un hilo
    de color de $R=10\Omega$. Si se conecta a $V=445V$ y la longitud
    del solenoide es de $l=20cm$. Calcular:
    \begin{itemize}
      \item Intensidad $I$.
      \item Inducción Magnética $B$.
    \end{itemize}
  \item Sabiendo que la inducción de un campo magnético uniforme es
    de $B=1,2T$ tesla. Calcular el flujo magnético $\Phi$ en un
    cuadrado de $L=0,5cm$ de lado perpendicular a las lineas de
    fuerza del campo magnético.
  \item Un solenoide de longitud $l=30cm$ y de radio $r=2cm$ está
    formado por $N=200$ espiras y es recorrido por una intensidad de
    $I=1A$. Calcular la inducción magnética $B$ en el interior del
    solenoide, y el flujo en Weber $Wb$ y Maxwell $Mx$.
  \item La inducción de un campo magnético uniforme es de
    $B=10000Gauss$. Calcular el flujo magnético que atraviesa la
    espira circular de radio $r=2cm$, colocada perpendicularmente a
    las líneas de fuerza.
  \item Por un solenoide de $N=180$ espiras, longitud $l=30cm$ y
    diámetro de $d=2cm$ circulan $I=2A$ amperes. Si están bobinados
    en un carrete de un núcleo de aire. Calcular expresando los
    resultados en unidades de los sistemas \emph{CGS y MKS(SI)
    sistema internacional}.
    \begin{center}
      \begin{tabular}{|c|c|} \hline
        CGS & MKS(SI) \\ \hline \hline
        Gauss [Gs] & Tesla [T] \\ \hline
        Maxwell [Mx] & Weber [Wb] \\ \hline
      \end{tabular}
    \end{center}
    \begin{enumerate}
      \item La inducción magnética $B$ en el núcleo del solenoide.
      \item El flujo magnético $\Phi$ en el núcleo.
    \end{enumerate}
\end{enumerate}
\end{document}

